\section{The ordinary heat-seeded iteration near \texorpdfstring{$W$}{W}}
\label{main:section}

We give the analytic argument and specify the finite certificate on which its
numerical constants depend. The certificate evaluates known Fourier and
exponential-polynomial expressions; it does not approximate an unknown
Navier--Stokes trajectory. Its exact rational data and checkers form part of
the accompanying reproducibility material.

\subsection{Conventions and statement}
\label{main:conventions}

Let $\mathbb T^3=(\mathbb R/2\pi\mathbb Z)^3$, with measure
$(2\pi)^{-3}\,dx$. Fourier coefficients are denoted by $\widehat f_k$;
$|\widehat f_k|$ is the Hermitian Euclidean norm on $\mathbb C^3$.
All vector fields are real, mean-zero and solenoidal. Thus
$\widehat f_{-k}=\overline{\widehat f_k}$ and
$k\cdot\widehat f_k=0$, where the latter product is bilinear.
For $k\ne0$, put $P_k=I-k\otimes k/|k|^2$, and write $\mathbb P$ for the
Leray projection. Define
\begin{align}
 A_s(f)&=\sum_{k\ne0}|k|^s|\widehat f_k|,
 & S(t)&=e^{t\Delta}, & A&=-\Delta,\label{main:norms}\\
 B(f,g)&=\mathbb P((f\cdot\nabla)g),
 &B_s(f,g)&=\tfrac12\{B(f,g)+B(g,f)\},\label{main:B}\\
 Q(f,g)(t)&=\int_0^t S(t-s)B_s(f,g)(s)\,ds,
 &\mathcal H r(t)&=\int_0^t S(t-s)r(s)\,ds.\label{main:Q}
\end{align}
The viscosity is one and there is no external force. The projected equation
is $u'+Au+B(u,u)=0$, and its mild form is $u=Su_0-Q(u,u)$.
In particular the sign convention in \eqref{main:B} is positive convection.
Let
\begin{equation}
 \mathcal X(I)=L^2(I;A_0),\qquad
 \|f\|_{\mathcal X(I)}^2=\int_I A_0(f(t))^2\,dt,
 \qquad \mathcal X=\mathcal X((0,\infty)).\label{main:X}
\end{equation}
Precisely, this is the closed real, mean-zero, solenoidal subspace of
$L^2(I;\ell^1(\mathbb Z^3\setminus\{0\};\mathbb C^3))$.

For $N>0$, $P_N$ retains the modes $0<|k|\le N$. The Galerkin solution
associated with $u_0$ solves
\begin{equation}
 u_N'+Au_N+P_NB(u_N,u_N)=0,\qquad u_N(0)=P_Nu_0.\label{main:galerkin}
\end{equation}
Write
\begin{equation}
 E_0=\|u_0\|_2^2,\quad X_0=\|\nabla u_0\|_2^2,\qquad
 B_*(u_0)=\sup_{N>0}\int_0^\infty\|\nabla u_N(t)\|_2^4\,dt.
 \label{main:Bstar}
\end{equation}

\begin{theorem}[Ordinary heat-seeded convergence]
\label{main:theorem}
Let $W=(\sin z,\sin x,\sin y)$ and $u_0=W+h$, where $h$ is smooth, real,
mean-zero and divergence-free, and
\begin{equation}
 A_{-1}(h)\le\varepsilon:=10^{-4}.\label{main:radius}
\end{equation}
There is a unique global smooth solution $u_h$ with datum $u_0$. The ordinary
Picard sequence
\begin{equation}
 v_0^h=S(W+h),\qquad v_{n+1}^h=T_h(v_n^h),\qquad
 T_h(f)=S(W+h)-Q(f,f)\label{main:picard}
\end{equation}
converges to $u_h$ in $\mathcal X$ on the entire half-line, with
\begin{equation}
 \|v_n^h-u_h\|_{\mathcal X}\le6(9/10)^{n-3}\quad(n\ge3),
 \qquad \int_0^\infty A_0(u_h)^2\,dt<9.\label{main:headline}
\end{equation}
For two perturbations satisfying \eqref{main:radius},
\begin{equation}
 \|u_h-u_g\|_{\mathcal X}\le4757A_{-1}(h-g).\label{main:stability}
\end{equation}
Moreover,
\begin{equation}
 B_*(u_0)\le\frac{E_0X_0}{2}e^9.\label{main:headline-Bstar}
\end{equation}
There is no restriction on the Fourier support of $h$.
\end{theorem}

The convergence rate concerns $\mathcal X$, not an $H^1$ rate. The quantities
$E_0$ and $X_0$ in \eqref{main:headline-Bstar} belong to the full datum; the
critical ball \eqref{main:radius} does not bound them uniformly. The theorem
is distinct from the larger regularity regions proved elsewhere in this
paper: their existence conclusions do not imply convergence of
\eqref{main:picard} throughout those regions.

\subsection{The heat bilinear estimate and the initial trace}
\label{main:heat-space}

\begin{lemma}[Complex pair bound and heat estimate]
\label{main:heat-lemma}
The space $\mathcal X$ is complete, and
\begin{equation}
 \|Q(f,g)\|_{\mathcal X}\le c\|f\|_{\mathcal X}\|g\|_{\mathcal X},
 \qquad c=\frac1{2\sqrt2}.\label{main:bilinear}
\end{equation}
The same bound holds for output projections and for zero-extended temporal
block inputs. For $A_{-1}(a)<\infty$,
\begin{equation}
 \|Sa\|_{\mathcal X}\le\frac{A_{-1}(a)}{\sqrt2}.\label{main:heat-data}
\end{equation}
The operator $Q$ maps $\mathcal X\times\mathcal X$ continuously into
$C([0,\infty);A_{-1})$, with zero initial trace, and
\begin{equation}
 \sup_{t\ge0}A_{-1}(Q(f,g)(t))
 \le\tfrac12\|f\|_{\mathcal X}\|g\|_{\mathcal X}.\label{main:trace}
\end{equation}
\end{lemma}

\begin{proof}
For a real unit vector $n$, write $a=\alpha n+a_\perp$ and
$b=\beta n+b_\perp$, allowing complex components. Then
\begin{align*}
 |P_{n^\perp}[(n\cdot a)b+(n\cdot b)a]|
 &=|\alpha b_\perp+\beta a_\perp|\\
 &\le |\alpha||b_\perp|+|\beta||a_\perp|
 \le |a||b|.
\end{align*}
The last inequality is Cauchy--Schwarz for the two nonnegative pairs.
Using input solenoidality to place derivatives on the output, and
reindexing the second symmetric convolution, gives
\begin{equation}
 |\widehat{B_s(f,g)}_k|
 \le\frac{|k|}{2}\sum_{p+q=k}|\widehat f_p||\widehat g_q|,
 \qquad k\ne0.\label{main:coefficient}
\end{equation}
The zero output coefficient vanishes. For each nonzero mode,
\[
 \big\|\mathbf1_{t\ge0}|k|e^{-|k|^2t}\big\|_{L^2_t}=1/\sqrt2.
\]
Minkowski in the Fourier sum, Young's $L^1_t$--$L^2_t$ convolution
inequality for this individual kernel, and time Cauchy--Schwarz prove
\eqref{main:bilinear}. This argument keeps the modewise kernel until its
time norm is taken; it does not integrate a pointwise supremum over modes.
For the homogeneous term, Minkowski gives
\[
 \|Sa\|_{\mathcal X}
 \le\sum_{k\ne0}|\widehat a_k|
       \bigg(\int_0^\infty e^{-2|k|^2t}\,dt\bigg)^{1/2}
 =\frac1{\sqrt2}A_{-1}(a).
\]
Dividing \eqref{main:coefficient} by $|k|$ and using contraction of $S$
in $A_{-1}$ proves
\[
 A_{-1}(Q(f,g)(t))\le\tfrac12\int_0^t A_0(f(s))A_0(g(s))\,ds,
\]
and hence \eqref{main:trace}. The forcing belongs to $L^1_tA_{-1}$.
Strong continuity of the heat semigroup on the weighted $\ell^1$ space,
together with approximation of that forcing by simple functions, proves
continuity and the zero trace. The same reasoning gives the trace $a$ for
$Sa$. One may first make all computations on finite Fourier sums and
compact time supports, then pass by density. Completeness follows from the
Banach-space definition of $\mathcal X$.
\end{proof}

We will also use uniqueness among mild solutions in $\mathcal X$ with the
same datum. Indeed, subdivide time into finitely many intervals, including
a final infinite interval, so that
$c(\|u\|_{\mathcal X(I)}+\|v\|_{\mathcal X(I)})<1$ on each one.
Their finite total squared norms permit such a subdivision. If the two
solutions agree on the preceding intervals, the entire preceding Duhamel
history cancels in their difference. Equation \eqref{main:bilinear} on the
next interval then forces equality there. This proves uniqueness on the
whole half-line without resetting either solution.

\subsection{Explicit references and the complete heat defect}
\label{main:reference}

For a smooth reference $w$ with initial trace $W$, define
$r_w=w'+Aw+B_s(w,w)$. Direct integration gives
\begin{equation}
 w-\mathcal H r_w=SW-Q(w,w).\label{main:correction-map}
\end{equation}
Thus, with $z=\mathcal H r_w$ and $\widetilde w=w-z$,
\begin{equation}
 r_{\widetilde w}=-2B_s(w,z)+B_s(z,z),\qquad z(0)=0.
 \label{main:correction-residual}
\end{equation}
These are identities for the full residual, including its low modes.

Let $w_0=SW=e^{-t}W$, $w_{j+1}=SW-Q(w_j,w_j)$, and put
\begin{align}
 F&=B(W,W)=(\sin y\cos z,\sin z\cos x,\sin x\cos y),\nonumber\\
 C&=\tfrac12(\sin z\cos2y,\sin x\cos2z,\sin y\cos2x),\nonumber\\
 G&=B(W,F)+B(F,W)=-W/2+C,\qquad H=B(F,F).\label{main:fields}
\end{align}
The Leray projection is part of these definitions. In particular $F\ne0$.
Define the nonnegative causal heat profiles
\begin{align}
 L_1&=\tfrac14\{e^{-t}-(1+2t)e^{-3t}\},\nonumber\\
 L_5&=\tfrac14\{e^{-5t}+(2t-1)e^{-3t}\},\nonumber\\
 L_6&=\tfrac14\{(2t^2-2t+1)e^{-4t}-e^{-6t}\}.\label{main:L}
\end{align}
They vanish at zero and satisfy
$(\partial_t+1)L_1=(\partial_t+5)L_5=te^{-3t}$ and
$(\partial_t+6)L_6=t^2e^{-4t}$. Consequently
\begin{align}
 w_1&=e^{-t}W-te^{-2t}F,\nonumber\\
 w_2&=(e^{-t}-L_1/2)W-te^{-2t}F+L_5C-L_6H.\label{main:w2}
\end{align}
The support of $w_2$ lies in $|k|^2\le6$. Its exact amplitude is
\begin{equation}
 U_2=A_0(w_2)=\tfrac{21}{8}e^{-t}+3te^{-2t}
                 +\tfrac32te^{-3t}+\tfrac38e^{-5t}+\sqrt3L_6.
 \label{main:U2}
\end{equation}
This follows from its disjoint Fourier shells and the positive fundamental
coefficient; it is not the sum of magnitudes of an arbitrarily signed
representation.

For explicit residual linkage, set $x=e^{-t}$, $y=te^{-2t}$,
$p=L_1/2$, $q=L_5$, $r=L_6$, and $K_{AB}=B(A,B)+B(B,A)$.
The scalar profiles $p,q,r$ are local to this subsection; the later block
argument uses $p$ for a positive vector and $q$ for a contraction factor.
Then
\begin{align}
 r_2={}&(p^2-2xp)F+ypG+(x-p)qK_{WC}-(x-p)rK_{WH}\nonumber\\
       &-yqK_{FC}+yrK_{FH}+q^2B(C,C)-qrK_{CH}+r^2B(H,H).
       \label{main:r2}
\end{align}
Define
\begin{equation}
 z_3=\mathcal H r_2,\qquad w_3=w_2-z_3,\qquad
 z_4=\mathcal H\{-2B_s(w_2,z_3)+B_s(z_3,z_3)\}.
 \label{main:w3z4}
\end{equation}
Thus $w_3$ is the third unperturbed heat-seeded iterate and
$T_h(w_3)-w_3=Sh-z_4$. The finite certificate constructs all signed
coefficients of $z_3$ (342 nonzero modes); the conservative support bound
$|k|^2\le24$ contains $w_3$. It constructs the signed mixed contribution
$-2\mathcal H B_s(w_2,z_3)$ to $z_4$ exactly. The quadratic contribution
$\mathcal H B_s(z_3,z_3)$ is bounded as a whole, not discarded or represented
as part of the mixed-mode count.

The reference calculations admit a real rational storage convention
without restricting the perturbation class. Under simultaneous inversion
of the spatial coordinates, $W(-\mathbf{x})=-W(\mathbf{x})$. If $f$ and $g$ are odd, then
$\nabla g$ is even, so $(f\cdot\nabla)g$ is odd; the Leray projection
and heat semigroup preserve this parity. Thus all the $W$-based reference
and defect fields constructed here are real and odd. Reality and oddness
together imply
\[
 \widehat f_k=i\xi_k,\qquad \xi_k\in\mathbb R^3,
 \qquad \xi_{-k}=-\xi_k,
 \qquad |\widehat f_k|=|\xi_k|.
\]
For two such fields, the positive-convection convention
\eqref{main:B} reads
\[
 \widehat{B(f,g)}_k
 =-iP_k\sum_{p+q=k}(\xi_p\cdot q)\eta_q,
 \qquad \widehat g_q=i\eta_q.
\]
This explains the minus sign in the stored real-vector convolution.
The programs store $\xi_k$, or the rational coefficients of its real
exponential-polynomial profile; their real Euclidean norm is exactly the
Hermitian norm of the actual Fourier coefficient. This convention is used
only for the explicit reference and defect calculations. The perturbation
$h$, its Picard iterates, and the fixed-point space are not assumed odd;
their estimates use the general complex-vector inequality
\eqref{main:coefficient}.

Here is the analytic basis of the finite enclosures. Heat convolution of
$t^m e^{-\alpha t}$ at frequency $|k|^2=\lambda$ is
$e^{-\lambda t}\int_0^t s^m e^{(\lambda-\alpha)s}\,ds$, a rational
exponential polynomial for the integer rates occurring here. Products of
mode profiles therefore have exact factorial moments. Their tails use
\begin{equation}
 \int_a^\infty t^m e^{-\lambda t}\,dt
 =e^{-\lambda a}\sum_{j=0}^m
       \frac{m!\,a^j}{j!\lambda^{m-j+1}}.\label{main:moments}
\end{equation}
After combining signed terms within each mode, Minkowski bounds its
$\mathcal X(I)$ norm by the sum of square roots of these squared moments.
Endpoint exponentials and radicals are enclosed outward by rational
series and integer-square-root bounds.

For the quadratic part of $z_4$, we construct a pointwise upper profile
explicitly. Let $\Phi_1,\ldots,\Phi_9$ be the nine spatial fields in the
order displayed in \eqref{main:r2}. Then
$r_2=\sum_{j=1}^9\sigma_j g_j\Phi_j$, where
\begin{align}
 (\sigma_1,\ldots,\sigma_9)&=(-1,1,1,-1,-1,1,1,-1,1),\nonumber\\
 (g_1,\ldots,g_9)&=\bigl(p(2x-p),yp,(x-p)q,(x-p)r,
                  yq,yr,q^2,qr,r^2\bigr).
 \label{main:positive-forcing}
\end{align}
The heat-convolution definitions of $L_1,L_5,L_6$ give
$x,y,p,q,r\ge0$, and the explicit formula for $L_1$ gives
$0\le p=L_1/2\le x/8$. Thus every $g_j$ is nonnegative on
$[0,\infty)$. For every shell present in $\Phi_j$, choose a rational
number
\[
 \eta^+_{j\lambda}\ge
 \sum_{|k|^2=\lambda}|\widehat{\Phi_j}_k|,
 \qquad \lambda\ge1,
\]
by summing outward upper bounds for the individual square roots. Define
\begin{equation}
 \psi_{j\lambda}(t)=\int_0^t e^{-\lambda(t-s)}g_j(s)\,ds,
 \qquad
 Z(t)=\sum_{j,\lambda}\eta^+_{j\lambda}\psi_{j\lambda}(t).
 \label{main:positive-causal-profile}
\end{equation}
All these sums are finite. Positivity of the scalar heat kernel implies
$\psi_{j\lambda}\ge0$. Since
$\widehat z_{3,k}=\sum_j\sigma_j\widehat{\Phi_j}_k
\psi_{j,|k|^2}$, the triangle inequality, first over fields and then over
shells, gives $A_0(z_3(t))\le Z(t)$. No disjointness of the fields' Fourier
supports is assumed. This is the certificate's 31-term field-shell
profile, stored under \texttt{pointwise\_A0\_upper\_profile} in
\nolinkurl{third_heat_defect_static_result.json}. Its positivity follows
from \eqref{main:positive-causal-profile}, even though the expanded
exponential-polynomial coefficients can have either sign.
For any $[a,b]$, including
$b=\infty$, the same modewise kernel calculation as above gives
\begin{equation}
 \|\mathcal H B_s(z_3,z_3)\|_{\mathcal X(a,b)}
 \le c\left\{\int_0^a e^{-(a-s)}Z(s)^2\,ds
                 +\int_a^b Z(s)^2\,ds\right\}.\label{main:quadratic-history}
\end{equation}
In particular its earlier forcing is retained. This estimate uses only
$|k|\ge1$; the corrected residual need not stay at high frequencies.

\subsection{A finite-block criterion for the unchanged global map}
\label{main:blocks}

Use the 17 intervals $I_i=[a_i,b_i]$ with endpoints
\begin{equation}
 0,\tfrac1{16},\tfrac18,\tfrac3{16},\tfrac14,\tfrac5{16},
 \tfrac38,\tfrac7{16},\tfrac12,\tfrac58,\tfrac34,\tfrac78,
 1,\tfrac54,\tfrac32,2,3,\infty.\label{main:partition}
\end{equation}
Overlaps of endpoints have measure zero. For $j<i$, let $k_{ij}$ be a
rational upper enclosure of
\begin{equation}
 \sup_{\lambda\ge1}e^{-\lambda(a_i-b_j)}
          \sqrt{1-e^{-2\lambda(b_i-a_i)}};\qquad k_{ii}=1.
 \label{main:memory}
\end{equation}
For the infinite block the square-root factor is one. Define rational
$s_i$ satisfying
\begin{equation}
 s_i\ge\frac1{\sqrt2}\sup_{\lambda\ge1}e^{-\lambda a_i}
       \sqrt{1-e^{-2\lambda(b_i-a_i)}}.
 \label{main:heat-block}
\end{equation}
Put
\[
 \kappa_i:=\sup_{\lambda\ge1}e^{-\lambda a_i}
       \sqrt{1-e^{-2\lambda(b_i-a_i)}},
 \qquad \overline\kappa_i\in\mathbb Q,
 \qquad \kappa_i\le\overline\kappa_i\le1.
\]
For $b_i=\infty$, the square-root factor is understood to be one;
then $\kappa_i=e^{-a_i}$. In the first block $a_1=0$ and
$\kappa_1=\overline\kappa_1=1$. The function
\texttt{heat\_factor} in
\nolinkurl{perturbation_block_gate.py} constructs these rational upper
enclosures. The certificate takes $s_i=2\bar c\,\overline\kappa_i$,
which satisfies \eqref{main:heat-block} because
$2\bar c>1/\sqrt2$, with
\begin{equation}
 \bar c=\frac{707106781187}{2000000000000}>c.
 \label{main:cbar}
\end{equation}
The first supremum in \eqref{main:heat-block} is one, even for a very short
first block. Arbitrarily high-frequency perturbations prevent a uniform
square-root-of-duration gain there.

\begin{lemma}[Complete source-block memory]
\label{main:block-lemma}
For $f,g\in\mathcal X$,
\begin{equation}
 \|Q(f,g)\|_{\mathcal X(I_i)}
 \le c\sum_{j\le i}k_{ij}
       \|f\|_{\mathcal X(I_j)}\|g\|_{\mathcal X(I_j)}.
 \label{main:block-bilinear}
\end{equation}
Also $\|Sh\|_{\mathcal X(I_i)}\le s_i A_{-1}(h)$.
\end{lemma}
\begin{proof}
Split the forcing into its source intervals. For $s\in I_j$, $j<i$, and
$\lambda=|k|^2$, the exact output-block norm of the differentiated heat
kernel is
\[
 \bigg(\int_{a_i}^{b_i}|k|^2e^{-2|k|^2(t-s)}\,dt\bigg)^{1/2}
 =\frac1{\sqrt2}e^{-\lambda(a_i-s)}
       \sqrt{1-e^{-2\lambda(b_i-a_i)}}\le k_{ij}/\sqrt2.
\]
Minkowski, \eqref{main:coefficient}, and source-time Cauchy--Schwarz give
each summand of \eqref{main:block-bilinear}. On the diagonal, use the whole
future kernel norm $1/\sqrt2$. Blocks later than $i$ contribute nothing.
The homogeneous heat estimate follows by the same Fourier Minkowski
calculation with no derivative. Every earlier source block remains in the sum.
\end{proof}

The suprema in \eqref{main:memory}--\eqref{main:heat-block} have elementary
enclosures. For gap $g>0$ and length $\ell>0$, the unconstrained squared
maximum over $\lambda\ge0$ is
$\ell(g+\ell)^{-1}(g/(g+\ell))^{g/\ell}$, with maximizer
$\log(1+\ell/g)/(2\ell)$. If that maximizer is at most one, the
constrained maximum is attained at one. The zero-gap and infinite-length
cases are treated directly. Rational logarithm, exponential, and square-root
enclosures suffice; no floating-point maximization is part of acceptance.

Choose certified rational bounds
\begin{equation}
 V_i\ge\|w_3\|_{\mathcal X(I_i)},\qquad
 d_i\ge\|z_4\|_{\mathcal X(I_i)}.\label{main:Vd}
\end{equation}
Here $V_i$ can be the sum of the enclosed $w_2$ and $z_3$ block norms.
Suppose positive rational radii $\rho_i$ satisfy
\begin{equation}
 d_i+\varepsilon s_i+
 \bar c\sum_{j\le i}k_{ij}(2V_j\rho_j+\rho_j^2)\le\rho_i.
 \label{main:invariance}
\end{equation}
The set
\begin{equation}
 \mathcal B=\{f\in\mathcal X:
       \|f-w_3\|_{\mathcal X(I_i)}\le\rho_i\text{ for all }i\}
 \label{main:rectangle}
\end{equation}
is then invariant under every $T_h$ in \eqref{main:radius}. Indeed,
$T_h(w_3)-w_3=Sh-z_4$, and expansion around $w_3$ followed by
\eqref{main:block-bilinear} proves \eqref{main:invariance}'s assertion.
On this set, block distances are bounded by the nonnegative lower-triangular
matrix
\begin{equation}
 M_{ij}=2\bar c\,k_{ij}(V_j+\rho_j)\ (j\le i),
 \qquad M_{ij}=0\ (j>i).\label{main:M}
\end{equation}
If $\max_i M_{ii}<1$, choose $q$ between this maximum and one and set
$p=(qI-M)^{-1}\mathbf1$. Forward substitution gives $p>0$ and
$Mp=qp-\mathbf1\le qp$. The norm
$\max_i\|f\|_{\mathcal X(I_i)}/p_i$ is equivalent to $\mathcal X$
because there are finitely many blocks. The map $T_h$ is a strict contraction
of the closed set \eqref{main:rectangle} in that complete norm.
This is an all-time fixed point of the unchanged map, not a chain of restarted
local iterations. Lemma~\ref{main:heat-lemma} gives its trace $W+h$.

\subsection{The finite certificate and entry of the actual heat seed}
\label{main:certificate}

The invariant-set construction alone would only prove convergence from
points already in $\mathcal B$. We must check entry for the particular seed
in \eqref{main:picard}. Let $V_{n,i}$ bound
$\|w_n\|_{\mathcal X(I_i)}$ for $n=0,1,2$. The first two exact amplitudes
are $3e^{-t}$ and $3e^{-t}+3te^{-2t}$. Their squared tail integrals are
\begin{align}
 F_0(t)&=\tfrac92e^{-2t},\nonumber\\
 F_1(t)&=\tfrac92e^{-2t}+(6t+2)e^{-3t}
            +(\tfrac94t^2+\tfrac98t+\tfrac9{32})e^{-4t}.
 \label{main:prefix-moments}
\end{align}
The $n=2$ bounds use \eqref{main:U2}. Thus these prefix inputs have explicit
analytic constructions independent of the unknown solution.

Set $E_{0,i}=\varepsilon s_i$ and choose outward rational upper bounds
recursively so that, for $n=0,1,2$,
\begin{equation}
 E_{n+1,i}\ge\varepsilon s_i+
 \bar c\sum_{j\le i}k_{ij}(2V_{n,j}E_{n,j}+E_{n,j}^2).
 \label{main:prefix}
\end{equation}
Subtracting the exact perturbed and unperturbed Picard equations and applying
Lemma~\ref{main:block-lemma} proves inductively
\begin{equation}
 \|v_n^h-w_n\|_{\mathcal X(I_i)}\le E_{n,i},\qquad 0\le n\le3.
 \label{main:prefix-error}
\end{equation}
No earlier source time is removed in this recurrence.

\begin{proposition}[Finite exact certificate for the stated constants]
\label{main:finite-certificate}
For \eqref{main:partition} and \eqref{main:cbar}, the supplied rational
data, constructed from \eqref{main:fields}--\eqref{main:quadratic-history},
satisfy \eqref{main:memory}--\eqref{main:invariance} and
\begin{equation}
 \max_i M_{ii}<\tfrac35,\qquad
 E_{3,i}<\rho_i\quad(1\le i\le17),\qquad
 C_0:=\sum_i(V_i+\rho_i)^2<9.\label{main:accepted}
\end{equation}
With $E=E_3+\rho$, there is a rational vector $p>0$ satisfying
\begin{equation}
 p\ge E,\qquad Mp\le\tfrac9{10}p,\qquad \sum_i p_i^2<36.
 \label{main:rate-check}
\end{equation}
For $b=(I-M)^{-1}s$, the rational data also satisfy
\begin{equation}
 b_i=\frac{s_i+\sum_{j<i}M_{ij}b_j}{1-M_{ii}}>0,
 \qquad \sum_i b_i^2<4757^2.\label{main:stability-check}
\end{equation}
\end{proposition}

\begin{proof}[Certificate construction and acceptance]
The following files specify the finite arithmetic proof; stage-source paths
are relative to \texttt{proof\_sources/W\_research/} in the companion
\texttt{Code.zip}. The replay-wrapper path is relative to that archive's root.
\begin{enumerate}
\item \nolinkurl{12_second_correction/second_correction_fourier_check.py}
and its result construct the projected fields and full second residual.
\item \nolinkurl{13_iteration_convergence/third_heat_defect_static.py}
and its result construct every signed $z_3$ mode, check its forced heat
identity and initial trace, and supply the positive profile $Z$.
\item \nolinkurl{14_unit_W_picard/fourth_heat_defect_static.py}
and its result construct the mixed $z_4$ modes and include the entire
quadratic contribution through \eqref{main:quadratic-history}.
\item \nolinkurl{14_unit_W_picard/refined_W_block_gate_check.py}
and \nolinkurl{refined_W_block_gate_result.json} evaluate the block moment enclosures and complete memory
factors. The replay wrapper \nolinkurl{checks/replay_W_chain.py} places
regenerated dependencies in a common \nolinkurl{working_copy/} directory,
where the Stage15 programs consume them using their relative paths.
\item \nolinkurl{15_perturbed_picard/perturbation_block_gate.py}
and \nolinkurl{perturbation_block_gate_result.json} provide the rational
$V,d,k,s,\rho,M$ and check every inequality \eqref{main:invariance}.
\item \nolinkurl{15_perturbed_picard/perturbed_picard_prefix_static.py}
and its result provide all $E_0,E_1,E_2,E_3$, verify
\eqref{main:prefix} with outward rounding, and check all 17 strict entry
comparisons in \eqref{main:accepted}.
\item \nolinkurl{15_perturbed_picard/perturbed_stability_rate_static.py}
and its result check the matrix formula, \eqref{main:rate-check},
\eqref{main:stability-check}, and the squared solution budget.
\end{enumerate}
Where not explicitly named, ``its result'' means the correspondingly named
\nolinkurl{_result.json} file. All integer/rational comparisons are performed
on the stored fractions; displayed decimals do not decide acceptance.
The exponential-polynomial integrals use \eqref{main:moments}. Endpoint
exponentials, logarithms used in memory factors, and radicals have explicit
outward series bounds. Source hashes bind the imported rational data; the
separate replay record identifies which generators and checkers were rerun.

For orientation only, the resulting values are
\begin{align*}
 \max_i M_{ii}&=0.5781572417201134\ldots,\\
 \max_i E_{3,i}/\rho_i&=0.9127337584808773\ldots,\\
 C_0&=8.340335743815434\ldots.
\end{align*}
The rate vector may be constructed by forward substitution:
\[
 p_i=\max\left\{E_i,
       \frac{\sum_{j<i}M_{ij}p_j}{9/10-M_{ii}}\right\}.
\]
Its domination and matrix inequalities follow directly from this formula;
the finite squared-norm comparison in \eqref{main:rate-check} is included
in the exact certificate. Thus the proposition rests on specified finite
identities, enclosures and inequalities, rather than a saved status string.
It remains a computer-assisted finite certificate, not a proof-assistant
formalization; the full rational arrays are supplied to permit inspection
and reproduction.
\end{proof}

By \eqref{main:prefix-error} and Proposition~\ref{main:finite-certificate},
the actual $v_3^h$ belongs to $\mathcal B$. If $u_h$ is its fixed point,
its block error after entry obeys
\begin{equation}
 \big(\|v_n^h-u_h\|_{\mathcal X(I_i)}\big)_i
 \le M^{n-3}(E_3+\rho),\qquad n\ge3.\label{main:vector-rate}
\end{equation}
The $E_3$ term is required because the entering iterate is perturbed while
the rectangle is centered at unperturbed $w_3$. Positivity and
\eqref{main:rate-check} give $M^{n-3}(E_3+\rho)\le(9/10)^{n-3}p$.
The block intervals partition the half-line, so their squared norms add.
This proves the rate in \eqref{main:headline}. Likewise,
\begin{equation}
 \|u_h\|_{\mathcal X}^2\le\sum_i(V_i+\rho_i)^2=C_0<9.
 \label{main:solution-budget}
\end{equation}

For $h,g$ in the same ball, let
$z_i=\|u_h-u_g\|_{\mathcal X(I_i)}$ and $\delta=A_{-1}(h-g)$.
Subtracting their fixed-point equations gives $z\le s\delta+Mz$.
Since $M_{ii}<1$, forward substitution yields $z\le b\delta$ with
$b$ as in \eqref{main:stability-check}. Taking the Euclidean norm proves
\eqref{main:stability}. Both endpoints must lie in the certified ball;
no assumption that their difference has norm at most $\varepsilon$ is used.

\subsection{Uniform Galerkin bounds and the full equation}
\label{main:pde}

We finish the regularity and all-cutoff claims without presupposing global
smoothness of the full equation. For $N\ge\sqrt{24}$, $P_Nw_3=w_3$.
On the Galerkin subspace the map
\[
 T_{h,N}(f)=SP_N(W+h)-P_NQ(f,f)
\]
has center defect $P_N(Sh-z_4)$. Output projection contracts all Fourier
amplitude norms. The same rectangle, invariant radii and Lipschitz matrix
therefore construct its global fixed point, with
\begin{equation}
 \int_0^\infty A_0(u_{h,N})^2\,dt\le C_0<9.
 \label{main:galerkin-budget}
\end{equation}
For completeness, this finite-dimensional mild fixed point is the
classical Galerkin solution. On every finite interval $[0,T]$, equivalence
of the finitely many mode norms gives
$P_NB(u_{h,N},u_{h,N})\in L^1(0,T)$ from
$u_{h,N}\in\mathcal X(0,T)$. Variation of constants therefore gives an
absolutely continuous representative satisfying
\[
 u_{h,N}'=-Au_{h,N}-P_NB(u_{h,N},u_{h,N})
 \quad\hbox{a.e.},\qquad
 u_{h,N}(0)=P_N(W+h).
\]
This representative is continuous, so the polynomial right-hand side is
continuous; hence the representative is $C^1$ and satisfies the ODE
everywhere. Local uniqueness for the finite-dimensional polynomial ODE
identifies it with the usual Galerkin solution. That solution extends
globally by its energy identity, since bounded energy prevents escape to
infinity in a fixed finite-dimensional space. This identification uses no
global regularity premise for the full PDE.

Put $X_N=\|\nabla u_{h,N}\|_2^2$ and
$Y_N=\|\Delta u_{h,N}\|_2^2$. Direct testing gives
\begin{align}
 \tfrac12X_N'+Y_N
 &\le\|u_{h,N}\|_\infty\sqrt{X_NY_N}
 \le\tfrac12Y_N+\tfrac12 A_0(u_{h,N})^2X_N,\nonumber\\
 X_N'+Y_N&\le A_0(u_{h,N})^2X_N.\label{main:H1}
\end{align}
The integrating factor and \eqref{main:galerkin-budget} imply
\begin{equation}
 \sup_{t\ge0}X_N(t)\le X_0e^{C_0},\qquad
 \int_0^\infty Y_N(t)\,dt\le X_0e^{C_0}.
 \label{main:strong-bounds}
\end{equation}
The energy equality gives $\int_0^\infty X_N\,dt\le E_0/2$, hence
\begin{equation}
 \int_0^\infty X_N(t)^2\,dt
 \le\frac{E_0X_0}{2}e^{C_0}<\frac{E_0X_0}{2}e^9
 \quad\hbox{if }E_0X_0>0.\label{main:large-cutoff}
\end{equation}
The zero datum has zero integral.

For $N<\sqrt{24}$ the reference need not belong to the Galerkin space;
\eqref{main:galerkin-budget} is not asserted. Energy instead gives
\[
 \sup_tX_N(t)\le N^2E_0<24E_0,\qquad
 \int_0^\infty X_N(t)^2\,dt\le12E_0^2\le12E_0X_0.
\]
The last step is the mean-zero Poincare inequality. Since $e^9>24$, these
cutoffs also satisfy \eqref{main:headline-Bstar}. This treatment does not
bound the full initial energy by the critical perturbation norm.

For completeness, \eqref{main:strong-bounds} provides uniform
$L^\infty_tH^1\cap L^2_tH^2$ bounds on every finite time interval.
The equation then bounds time derivatives in the usual spaces needed for
Galerkin compactness (one may use $L^2_tL^2_x$ here, since
$H^2(\mathbb T^3)\hookrightarrow L^\infty$ and
$\nabla u_N\in L^\infty_tL^2_x$).
Simon~\cite{Simon}, Corollary~4, applied with $H^2\Subset H^1\hookrightarrow L^2$
on each finite time interval, gives strong $L^2_tH^1_x$ compactness.
It passes the quadratic term to a global strong solution with the prescribed datum.
The standard local smooth solution agrees with this limit by strong
uniqueness and continues for all time under the uniform $H^1$ bound;
smooth initial data give smoothness at every finite time. The local $H^1$
existence, uniqueness and smoothing statements used here are reviewed in
Tao~\cite{Tao2009}, Proposition~2.2(i)--(ii). Its unit-torus convention
transfers to our fixed torus by spatial, temporal and amplitude rescaling;
we use only the qualitative result. Restarting its uniform local lifespan
under the bounded $H^1$ norm gives continuation. Applying Fatou
to the Fourier amplitudes and then time preserves the bound $C_0$.
The limit is therefore a mild solution in $\mathcal X$ and agrees with
the fixed point constructed above by the $\mathcal X$ uniqueness argument.
This proves all assertions of Theorem~\ref{main:theorem}.

\paragraph{Scope of the finite certificate.}
The temporal partition and all reference fields are finite, whereas the
perturbation and limiting equation are not truncated. The finite calculations
prove the bounds used by the infinite-dimensional analytic argument.
They do not establish success of repeated residual correction for arbitrary
data, nor do failed larger-radius sufficient tests establish divergence or
singularity. The exact arithmetic and its reproducibility record also do
not constitute external mathematical peer review or a novelty claim.
